\begin{table}[!h]
\centering
\caption{\label{tab:p8_inf}Respuesta de expectativas de inflación al shock cambiario}
\centering
\begin{threeparttable}
\begin{tabular}[t]{cccc}
\toprule
$h$ & Expectativa & Realizado 12m fwd & Gap (exp $-$ real)\\
\midrule
0 & -0.004 [-0.022, 0.014] & 0.075 [0.001, 0.148] & -0.082 [-0.164, 0.000]\\
3 & 0.022 [-0.002, 0.047] & 0.080 [0.007, 0.152] & -0.053 [-0.128, 0.022]\\
6 & 0.040 [-0.012, 0.092] & 0.051 [0.001, 0.100] & -0.005 [-0.075, 0.064]\\
12 & 0.009 [-0.020, 0.038] & 0.051 [-0.026, 0.128] & -0.039 [-0.132, 0.053]\\
18 & 0.018 [-0.003, 0.040] & 0.043 [-0.022, 0.108] & -0.024 [-0.082, 0.033]\\
\addlinespace
24 & 0.027 [-0.003, 0.056] & 0.019 [-0.018, 0.056] & 0.008 [-0.034, 0.049]\\
\bottomrule
\end{tabular}
\begin{tablenotes}
\item \textit{Note: } 
\item Coeficiente $beta_h$ de la LP con el shock cambiario $hat u_t$ del VAR bivariado (Secci'on 3) como regresor; controles: rezagos del VAR. emph{Expectativa}: respuesta de $pi^{e,12}_{t+h}$ (nivel). emph{Realizado 12m fwd}: respuesta de $IPC_{t+h+12}-IPC_{t+h}$, la inflaci'on efectivamente realizada en los 12 meses posteriores a $t+h$, comparable con la expectativa. emph{Gap}: diferencia entre ambas; su coeficiente es el test directo de igualdad de las dos respuestas. Intervalos al 95% con errores de Newey-West y $h+1$ rezagos. Ambas series en puntos log de inflaci'on acumulada a 12 meses. Muestra: 2001m12--2025m12.
\end{tablenotes}
\end{threeparttable}
\end{table}